% Source: https://artofproblemsolving.com/community/c3h1969377_question
% Stable identity: aops-topic-1969377-q01
For every real number $t$, $\sqrt{t^2}=|t|$, because the square root denotes the \emph{nonnegative} number whose square is $t^2$. So the equation says
\[
|a-2|=2-a=-(a-2).
\]
Now $|t|=-t$ exactly when $t\le0$. With $t=a-2$ this means $a-2\le0$. Hence the identity holds exactly for
\[
\boxed{a\le2}.
\]
It is not true for all $a$: for example, $a=3$ gives $\sqrt{1}=1$ on the left but $2-3=-1$ on the right.
